\begin{abstract}
Finding shortest paths from a single source to $k$ goals is well studied
in the algorithms literature but has received comparatively little attention
in heuristic search. $k$A$^*$ addresses this by performing a single search
toward all $k$ goals, aggregating heuristic values to guide exploration.
However, every generated node must be evaluated against all remaining goals,
making aggregation a bottleneck as $k$ grows. We present
iterative-$k$A$^*$, which solves goals sequentially while reusing a shared
open list and previously computed $g$-values, reducing heuristic computations
by up to a factor of $k$. We prove that, up to tie-breaking,
iterative-$k$A$^*$ expands the same set of nodes as
aggregative-$k$A$^*$. We further introduce a backward-search framework that
reverses the problem into many-to-one shortest paths, where each
goal-to-source search yields exact suffix costs that serve as
cached perfect heuristics for subsequent searches. When exact costs are
unavailable, the framework derives lower bounds from previously solved goals
and propagates them to tighten guidance and reduce expansions.
Experiments on 2,500 Moving AI 4-connected grid instances with Manhattan
heuristics show that iterative-$k$A$^*$ matches aggregative-$k$A$^*$ in
node expansions while achieving the lowest runtime among all variants. The
backward framework helps on some instances, though forward search dominates
overall - increasingly so as $k$ grows.
\end{abstract}
